\section{Pathologies and Limits}


本節では、本稿の枠組みがどこまで有効であり、どこから先を扱わないのかを整理する。

ここでいう pathologies とは、医学的診断名を意味しない。

本稿は、精神疾患、神経疾患、意識障害を診断する理論ではない。

ここで扱うのは、前節までに導入した consciousness window から外れたとき、ログ参照状態がどの方向へ崩れうるかという、構造的な逸脱である。

したがって、本節の目的は病理の説明ではなく、意識状態の限界条件を明確にすることである。

\subsection{Pathology as Regime Deviation}

本稿の枠組みでは、安定した意識状態は次のような窓として表される。

\begin{equation}
0 < |\Delta v| < \Delta v_{\max}
\end{equation}

\begin{equation}
\mu_{\min} < \mu < \mu_{\max}
\end{equation}

\begin{equation}
\tau_{\min} < \tau < \tau_{\max}
\end{equation}

\begin{equation}
\mathcal{T}_{\min}
<
\mathcal{T}_{\Phi}
<
\mathcal{T}_{\max}
\end{equation}

この窓の外側では、意識が単に「弱くなる」のではない。

崩れ方には方向がある。

速度差が消える方向。

速度差が過剰になる方向。

結合粘性が低すぎる方向。

結合粘性が高すぎる方向。

意味勾配が弱すぎる方向。

意味勾配が強すぎる方向。

測定や報告によって、ライブな境界が過剰に沈殿する方向。

それぞれが異なる参照状態の変形をもたらす。

ここで重要なのは、これらを疾患名に直結させないことである。

本稿が扱うのは、診断分類ではなく、時間スケール間のレジーム遷移である。

\subsection{Collapse Toward Synchronization}

第一の逸脱は、速度差が小さくなりすぎる方向である。

\begin{equation}
\Delta v \rightarrow 0
\end{equation}

このとき、高速層と低速層の境界は同期に潰れる。

処理は安定し、効率化されるが、参照のための時間的厚みは失われる。

この方向では、行動は自動化し、反応は滑らかになり、明示的な意識的参照は弱くなる。

これは必ずしも悪い状態ではない。

熟練運動、習慣化された行為、反射的処理の多くは、この方向に近い。

しかし、過度に同期側へ寄ると、状況の再解釈やログの再参照が起きにくくなる。

この場合、意識は停止するというより、不要化される。

つまり、処理は走るが、ログ参照として開かれる必要が薄くなる。

\subsection{Dissociation by Excessive Difference}

第二の逸脱は、速度差が大きくなりすぎる方向である。

\begin{equation}
|\Delta v| \rightarrow \Delta v_{\max}
\quad \text{or beyond}
\end{equation}

このとき、高速層と低速層の対応は不安定になる。

高速層では連想、予測、イメージ、反応が走るが、それらが低速層へ安定して沈殿しにくくなる。

この方向では、意識状態は断片化しやすい。

夢、白昼夢、没入、解離的状態、幻覚様状態は、このレジームの一部として記述しうる。

ただし、ここでも注意が必要である。

本稿は、これらの状態を同一視しない。

共通しているのは、速度差が大きくなり、ログ参照の安定性が低下するという構造だけである。

したがって、この節は診断分類ではなく、参照構造の崩れ方を示すものである。

\subsection{Fixation by Low Viscosity}

第三の逸脱は、結合粘性 $\mu$ が低すぎる場合である。

\begin{equation}
\mu \rightarrow 0
\end{equation}

この状態では、高速層と低速層が過剰に密着する。

速度差は十分に滑らず、処理は特定のログや意味井戸に強く結びつきやすくなる。

この方向では、反復思考、過集中、固着、強迫的参照のような状態が生じやすい。

重要なのは、ここで起きているのが「意味の強さ」だけではないという点である。

問題は、意味勾配と低粘性の結合によって、処理が同じ井戸へ何度も落ち込みやすくなることである。

すなわち、固定化は内容の問題ではなく、参照経路の粘性問題として扱われる。

この場合、意識内容は明瞭に感じられることがある。

しかし、その明瞭さは自由度の高さを意味しない。

むしろ、意味井戸への固定が強まり、再解釈可能性が低下している場合がある。

\subsection{Fragmentation by High Viscosity}

第四の逸脱は、結合粘性 $\mu$ が高すぎる場合である。

\begin{equation}
\mu \rightarrow \infty
\end{equation}

この状態では、高速層と低速層は互いに滑りすぎる。

高速な処理は走っていても、それが低速なログとして十分に沈殿しない。

その結果、経験は残りにくく、自己連続性や時間的まとまりは弱くなる。

この方向では、参照の断片化、記憶への沈殿不全、報告の不安定化が起こりうる。

ただし、本稿はこれを特定の疾患に対応づけない。

ここで述べているのは、あくまで結合粘性が高すぎる場合に予測される構造的傾向である。

処理は存在しても、ログ参照として安定しない。

この点が、同期側への崩れとは異なる。

\subsection{Weak Meaning Gradient}

第五の逸脱は、意味勾配が弱すぎる場合である。

\begin{equation}
|\nabla \Phi| \rightarrow 0
\end{equation}

この状態では、ログ参照は方向を失う。

処理は存在していても、どのログが参照されるべきか、どの概念井戸へ落ち込むべきかが定まりにくい。

この方向では、意識はぼんやりし、注意は散漫になり、経験の輪郭が弱くなる。

ただし、意味勾配が弱いことは、必ずしも機能低下を意味しない。

休息、拡散的思考、自由連想、創造的漂流は、この領域を利用する場合がある。

問題になるのは、意味勾配が弱すぎて、ログ参照が安定した地形を形成できない場合である。

つまり、弱い意味勾配は、柔らかい探索にもなりうるし、参照の喪失にもなりうる。

その違いは、速度差、粘性、遅延、ゲート状態との組み合わせによって決まる。

\subsection{Excessive Meaning Gradient}

第六の逸脱は、意味勾配が強すぎる場合である。

\begin{equation}
|\nabla \Phi| \rightarrow \infty
\end{equation}

この状態では、処理は特定の意味井戸へ強く引き込まれる。

ログ参照は明瞭になるが、自由度は低下する。

この方向では、感情的に強い記憶、トラウマ的再参照、強迫的反復、物語への過剰同一化のような状態が記述されうる。

ここで重要なのは、意味勾配が意識を強めるだけではないという点である。

勾配は参照を方向づけるが、過剰な勾配は参照を固定する。

したがって、意味が強いことは、必ずしも理解が深いことを意味しない。

強すぎる意味場は、むしろ理解の可塑性を奪うことがある。

この点で、$\mathcal{T}_{\Phi}$ は両義的である。

それはログ参照を開く条件であると同時に、過剰になれば参照を固定する条件にもなる。

\subsection{Measurement-Induced Collapse}

第七の逸脱は、測定や報告によって、ライブな境界が過剰に沈殿する場合である。

第13節で述べたように、測定はクオリアそのものを保存する操作ではない。

測定は、ライブな境界状態を低速層へ沈殿させ、報告可能な残余へ変換する操作である。

したがって、測定要求が強すぎる場合、意識状態はライブな開口としてではなく、過剰に安定化されたログとして扱われやすくなる。

このとき、残るものは明瞭である。

言語化され、数値化され、比較可能になる。

しかし、その明瞭さはライブなクオリアそのものではない。

それは、クオリアが沈殿した後の residue の明瞭さである。

この逸脱は、実験設計や内省において重要である。

観察しようとするほど、対象は変形される。

説明しようとするほど、ライブな境界はログ化される。

したがって、本稿の枠組みでは、測定や内省もまた意識状態を変化させるレジーム操作として扱われる。

\subsection{Mislocalization of Consciousness}

第八の逸脱は、意識を局所基盤や沈殿残余と取り違える場合である。

第14節で述べたように、意識には神経基盤がある。

しかし、意識状態そのものは単一部位に局在しない。

局在化によって捉えられるのは、局所的な基盤、神経相関、報告可能なログ、測定後の残余である。

これらは重要である。

しかし、それらを意識そのものと同一視すると、ライブなログ参照レジームが見えなくなる。

この誤りは、意識の座を探す議論だけでなく、自己やクオリアを脳内の特定対象として探す議論にも現れる。

本稿の立場では、局所基盤を否定しない。

しかし、局所基盤は意識の必要条件でありうるが、意識そのものではない。

意識は、基盤が特定の時間的・意味的・非閉包的関係に入ったときに開かれる状態である。

\subsection{Limits of the Present Framework}

ここまでの記述は、意識状態の逸脱を構造的に整理するためのものである。

しかし、本稿には明確な限界がある。

第一に、本稿は医学的診断理論ではない。

ここで述べたレジームは、精神医学的分類や神経疾患の分類に直接対応しない。

第二に、$\mathcal{T}_{\Phi}$ の測定可能な定義は、まだ十分に確立されていない。

\begin{equation}
\mathcal{T}_{\Phi}
\sim
\kappa
\frac{\Delta v}{\mu}
|\nabla \Phi|
\end{equation}

という関係は、現時点では構造式であり、具体的な測定量ではない。

第三に、$\Phi$、$I$、$J$ の IFGT 的定義を、神経活動・言語報告・主観報告へどのように写像するかは未完である。

第四に、クオリアを境界現象として扱うことは、主観経験のすべてを説明するものではない。

それは、クオリア問題の位置を変えるものであって、主観性の全体を消去するものではない。

第五に、本稿は意識の統一方程式を与えない。

意識を閉じた方程式に固定しようとすると、意識を成立させている時間的非閉包性そのものが失われる。

したがって、本稿の理論的射程は、意識そのものの完全記述ではなく、意識が成立しうる構造条件の指定にある。

\subsection{Why the Limits Matter}

これらの限界は、本稿の弱点であると同時に、理論の安全装置でもある。

意識を単一変数として測定しようとすれば、理論はすぐに生成モデルへ戻ってしまう。

クオリアを保存物として扱えば、境界現象としての性質が失われる。

病理を診断名に直結させれば、構造的レジームの柔軟性が失われる。

意識を局所部位に置けば、ライブなログ参照ループが見えなくなる。

したがって、本稿はあえて複数の点で閉じない。

閉じないことは、曖昧さではない。

本稿が扱う対象そのものが、非閉包的に維持される時間構造だからである。

理論が完全閉包を急ぐと、対象である意識を失う。

この制約を明示することが、本節の目的である。

\subsection{Summary}

本稿における pathologies とは、疾患分類ではなく、consciousness window からの逸脱方向である。

\begin{itemize}
\item $\Delta v \rightarrow 0$：同期化、自動化、意識的参照の弱化。
\item $|\Delta v|$ の過大化：断片化、夢様状態、解離様状態、幻覚様状態。
\item $\mu \rightarrow 0$：低粘性による固着、反復、意味井戸への固定。
\item $\mu \rightarrow \infty$：高粘性による沈殿不全、参照の断片化。
\item $|\nabla \Phi| \rightarrow 0$：意味勾配の消失、輪郭の弱化。
\item $|\nabla \Phi|$ の過大化：意味井戸への過剰固定、可塑性の低下。
\item 測定要求の過剰化：ライブなクオリアの沈殿残余化。
\item 局在化の過剰化：局所基盤や残余を意識そのものと取り違えること。
\end{itemize}

これらはいずれも、意識の有無を二値的に判定するものではない。

意識状態がどの方向に変形し、どの条件で安定窓から外れるかを示すための、構造的な分類である。
